% Общий шаблон тренировочного варианта PhysEGE. Только условия заданий.
\documentclass[12pt,a4paper]{article}
\usepackage{fontspec}
\IfFontExistsTF{Times New Roman}{\setmainfont{Times New Roman}}{\setmainfont{Liberation Serif}}
\usepackage{polyglossia}
\setdefaultlanguage{russian}
\usepackage{amsmath,amssymb,graphicx,tabularx,array}
\usepackage{geometry,fancyhdr,needspace}
\geometry{left=22mm,right=19mm,top=22mm,bottom=21mm,headheight=15pt}
\setlength{\parindent}{9mm}
\setlength{\parskip}{0pt}
\linespread{1.08}
\pagestyle{fancy}
\fancyhf{}
\fancyhead[R]{\small Физика. %%VARIANT_TITLE%% --- \thepage}
\fancyfoot[C]{\small %%YEAR%% г.}
\renewcommand{\headrulewidth}{0pt}
\renewcommand{\footrulewidth}{0pt}
\newcommand{\TaskStart}[1]{\par\medskip\needspace{45mm}\noindent\llap{\fbox{\makebox[9mm][c]{\bfseries #1}}\hspace{3mm}}}
\newcommand{\AnswerLine}{\par\smallskip\noindent Ответ:~\rule{68mm}{0.4pt}.\par\smallskip}
\newcommand{\AnswerBoxes}[1]{\par\smallskip\noindent Ответ:~\begingroup\count0=0\loop\ifnum\count0<#1\fbox{\rule{0pt}{5mm}\makebox[6mm]{}}\advance\count0 by1\repeat\endgroup\par\smallskip}
\newcommand{\PairAnswer}{\par\smallskip\noindent Ответ:~(\rule{26mm}{0.4pt}\,$\pm$\,\rule{22mm}{0.4pt}).\par\smallskip}
\newcommand{\MatchingAnswer}{\par\smallskip\noindent Ответ:~\begin{tabular}{|c|c|}\hline А & Б\\\hline\rule{0pt}{7mm}\hspace{7mm} & \hspace{7mm}\\\hline\end{tabular}\par\smallskip}
\newcommand{\TaskFigure}[1]{\par\smallskip\begin{center}\includegraphics[width=.78\linewidth,height=63mm,keepaspectratio]{#1}\end{center}\smallskip}
\newcommand{\ExamNotice}[1]{\par\medskip\noindent\fbox{\begin{minipage}{.965\linewidth}\textbf{\itshape #1}\end{minipage}}\par\medskip}
\begin{document}
\begin{center}
{\large\bfseries %%VARIANT_TITLE%%}\par\bigskip
{\large\bfseries Инструкция по выполнению работы}
\end{center}
Для выполнения экзаменационной работы по физике отводится 3 часа 55 минут (235 минут). Работа состоит из двух частей, включающих 26 заданий.

Ответом к заданиям 1--20 является число или последовательность цифр. Ответ запишите в поле ответа в тексте работы, а затем перенесите в бланк ответов № 1. Единицы измерения физических величин писать не нужно.

Ответы к заданиям 21--26 включают подробное описание всего хода выполнения задания. В бланке ответов № 2 укажите номер задания и запишите его полное решение.

При вычислениях разрешается использовать непрограммируемый калькулятор. При выполнении заданий можно пользоваться черновиком. Записи в черновике не учитываются при оценивании работы.

После завершения работы проверьте, чтобы ответ на каждое задание был записан под правильным номером.

\begin{center}\bfseries\itshape Желаем успеха!\end{center}
\clearpage
%%REFERENCE%%
\clearpage
\begin{center}\large\bfseries Часть 1\end{center}
\ExamNotice{Ответами к заданиям 1--20 являются число или последовательность цифр. Запишите ответ в поле после каждого задания.}
%%PART_ONE%%
\clearpage
\begin{center}\large\bfseries Часть 2\end{center}
\ExamNotice{Для записи ответов на задания 21--26 используйте бланк ответов № 2. Сначала укажите номер задания, затем запишите решение.}
%%PART_TWO%%
\end{document}
